% ==========================================
% BAB VII KESIMPULAN DAN SARAN
% ==========================================
\cleardoublepage
\chapter{KESIMPULAN DAN SARAN}
\label{chap:kesimpulan-saran}

\section{Kesimpulan}
Berdasarkan seluruh tahapan pengerjaan tugas akhir, mulai dari analisis kebutuhan, perancangan, implementasi, 
hingga evaluasi, dapat ditarik beberapa kesimpulan sebagai berikut.

\begin{enumerate}
    \item \textit{Frontend dashboard} berbasis web untuk sistem otomasi dosis WTP ITB Jatinangor berhasil 
    diimplementasikan dan diintegrasikan dengan layanan \textit{backend} melalui REST API. \textit{Dashboard} 
    yang dibangun menyatukan informasi kualitas air (pH, kekeruhan, dan TDS), status kepatuhan parameter, level 
    tangki bahan kimia, status pompa dosis, alarm, serta riwayat operasional ke dalam satu antarmuka pemantauan 
    yang terpusat. Dengan demikian, persoalan tersebarnya informasi operasional pada sumber yang terpisah, 
    sebagaimana dirumuskan pada Bab~\ref{chap:pendahuluan}, telah terjawab.

    \item Antarmuka yang dikembangkan tidak hanya berfungsi sebagai media visualisasi, tetapi juga mendukung 
    aksi operasional operator melalui kendali antarmuka yang terproteksi, yaitu pengakuan (\textit{acknowledge}) 
    alarm, peninjauan riwayat dengan penyaringan rentang waktu, ekspor laporan dalam format CSV dan XLSX, serta 
    pengiriman nilai \textit{offset} dosis sementara. Seluruh interaksi tersebut hanya dapat diakses setelah 
    operator melalui proses autentikasi, baik menggunakan kredensial lokal maupun penyedia identitas pihak ketiga.

    \item Persoalan risiko pengambilan keputusan berdasarkan data yang tidak valid ditangani melalui penanganan 
    kondisi data yang eksplisit pada antarmuka. Sistem membedakan secara visual antara kondisi data sedang dimuat, 
    data kosong, kegagalan pengambilan data, dan data kedaluwarsa, dengan komponen peringatan yang muncul ketika 
    telemetri terbaru tidak tersedia atau sistem berada dalam kondisi \textit{offline}.

    \item Hasil pengujian fungsional terhadap 20 skenario menunjukkan bahwa seluruh skenario berhasil dilalui 
    dengan status \textit{passed}. Hasil tersebut memverifikasi bahwa kesembilan kebutuhan fungsional yang 
    didefinisikan pada Bab~\ref{chap:analisis-masalah} telah terpenuhi, mencakup autentikasi, pemantauan 
    \textit{dashboard}, analitik sensor, peninjauan riwayat dan ekspor laporan, penanganan alarm, serta interaksi 
    \textit{offset} pompa.

    \item Hasil pengujian usabilitas menggunakan instrumen \textit{System Usability Scale} (SUS) memperoleh skor 
    75, yang melampaui target minimal 70 dan berada pada kategori \textit{good} serta \textit{acceptable}. Hasil 
    ini mengindikasikan bahwa antarmuka yang dikembangkan dapat dipahami dan digunakan oleh operator sasaran, 
    meskipun dengan catatan bahwa pengujian hanya melibatkan satu responden sehingga belum dapat digeneralisasi.

    \item Kontribusi utama subsistem \textit{frontend} pada tugas akhir ini bukan pada penentuan keputusan 
    kualitas air maupun perhitungan dosis bahan kimia, melainkan pada pengorganisasian data yang disediakan 
    \textit{backend} menjadi alur kerja pemantauan yang dapat digunakan operator. Dengan memisahkan pemantauan 
    ringkas, analitik sensor, peninjauan riwayat, penanganan alarm, dan intervensi dosis ke dalam halaman yang 
    berbeda, antarmuka mengurangi kebutuhan operator untuk berpindah antar sumber informasi selama proses 
    supervisi berlangsung.
\end{enumerate}

\section{Saran}
Berdasarkan hasil evaluasi dan batasan yang telah diidentifikasi, terdapat beberapa saran yang dapat menjadi 
arah pengembangan lebih lanjut terhadap subsistem \textit{frontend}.

\begin{enumerate}
    \item \textbf{Penerapan komunikasi dua arah yang persisten.} Mekanisme pengambilan data saat ini masih 
    mengandalkan penarikan berkala (\textit{polling}) sehingga terdapat jeda antara kondisi aktual di lapangan 
    dan tampilan pada antarmuka. Penerapan kanal komunikasi dua arah yang persisten dapat mempersingkat jeda 
    tersebut, khususnya untuk kejadian yang membutuhkan respons segera seperti pemicuan alarm.

    \item \textbf{Dukungan operasi luring melalui antrean mutasi.} Aksi pengakuan alarm dan pengiriman 
    \textit{offset} dosis saat ini bergantung sepenuhnya pada ketersediaan layanan \textit{backend}. Penambahan 
    mekanisme antrean mutasi luring (\textit{offline mutation queue}) memungkinkan aksi operator tetap tersimpan 
    ketika koneksi terputus dan dikirimkan kembali secara otomatis setelah koneksi pulih. Kemampuan ini relevan 
    mengingat area instalasi berpotensi memiliki titik buta sinyal nirkabel.

    \item \textbf{Peningkatan kejelasan prioritas alarm.} Tingkat keparahan alarm dapat dibedakan secara visual 
    dengan lebih tegas, misalnya melalui perbedaan warna maupun ikon antara alarm berkategori \textit{warning} 
    dan \textit{critical}sehingga operator dapat menentukan prioritas penanganan dengan lebih cepat.

    \item \textbf{Penambahan umpan balik konfirmasi pada aksi penting.} Berdasarkan pola jawaban pada pengujian 
    usabilitas, aspek pemanduan pengguna masih menyisakan ruang perbaikan. Penambahan pesan konfirmasi yang lebih 
    eksplisit pada aksi kritis, seperti pengakuan alarm dan pengiriman \textit{offset} dosis, dapat membantu 
    operator memastikan bahwa aksinya telah berhasil diproses oleh sistem.

    \item \textbf{Perluasan cakupan pengujian.} Pengujian usabilitas sebaiknya melibatkan jumlah responden yang 
    lebih banyak agar hasilnya memiliki kekuatan generalisasi. Selain itu, pengujian lanjutan dapat diperluas 
    untuk mengevaluasi keandalan antarmuka pada kondisi jaringan nyata di lingkungan instalasi, serta mengukur 
    waktu respons operator dalam jangka panjang.

    \item \textbf{Pengayaan visualisasi pada halaman utama.} Halaman \textit{dashboard} dapat dilengkapi dengan 
    grafik tren ringkas sehingga operator dapat mengamati arah perubahan kondisi air secara langsung dari layar 
    pemantauan utama tanpa perlu berpindah ke halaman analitik sensor.
\end{enumerate}